\clearpage
\begingroup
\fancyhead[L]{\small THEORIE DER BEWEGUNG DES MONDES.}
\renewcommand{\headrulewidth}{0pt}
\SourceLeaf{638}
\par\noindent\hspace*{\parindent}Da in dem Werthe von \(\theta\) die Grösse \(M\) in keinen Theilen einer höhern
Ordnung, als die dritte, vorkommt, so sieht man leicht, dass die Reduction
III nur Glieder der sechsten Ordnung hinzufügen würde; daher schlechtweg
\(M=M^{\prime}\) zu schreiben ist. — Nachdem alle Reductionen vorgenommen sind,
findet sich endlich, dass man, um \(\beta\) zu erhalten, in \(\theta\) statt \(E,\lambda,M,a\) resp.
\(E^{\prime},\lambda^{\prime},M^{\prime},a^{\prime}\) schreiben und noch folgende Theile hinzufügen müsse:
\begin{center}
\makebox[.75\linewidth][s]{.\hfill.\hfill.\hfill.\hfill.\hfill.\hfill.\hfill.\hfill.\hfill.\hfill.\hfill.\hfill.\hfill.\hfill.\hfill.\hfill.\hfill.\hfill.\hfill.\hfill.\hfill.\hfill.\hfill.\hfill.\hfill.\hfill.\hfill.\hfill.\hfill.\hfill.\hfill.\hfill.\hfill.}\\[4pt]
\makebox[.75\linewidth][s]{.\hfill.\hfill.\hfill.\hfill.\hfill.\hfill.\hfill.\hfill.\hfill.\hfill.\hfill.\hfill.\hfill.\hfill.\hfill.\hfill.\hfill.\hfill.\hfill.\hfill.\hfill.\hfill.\hfill.\hfill.\hfill.\hfill.\hfill.\hfill.\hfill.\hfill.\hfill.\hfill.\hfill.}
\end{center}
\par\noindent\hspace*{\parindent}Die \textsc{Mayer}schen Tafeln haben noch etwas Eigenthümliches. Es wird
nemlich darin nicht die wahre Anomalie \(M^{\prime}\), sondern ein gleichsam Mittelding
zwischen wahrer und mittlerer Anomalie, \(M^{\prime\prime}\), gebraucht, und an den
Ort des mittlern Knoten wird eine Verbesserung angebracht, die von \(a\) abhängt
und dazu dient, die Theile, deren Argumente \(a+\lambda\) und \(2a+\lambda\) sind,
wegzuschaffen. Diesem zufolge hat man
\[
\begin{aligned}
M^{\prime}={}&M^{\prime\prime}-2e\sin M^{\prime\prime}
 +\tfrac54ee\sin2M^{\prime\prime}\\
&+\tfrac{11}{8}nn\sin2E^{\prime}-\tfrac14gg\sin2\lambda
\end{aligned}
\]
\[
\lambda^{\prime}=\lambda^{\prime\prime}
 +\left(\tfrac94n\varepsilon-\tfrac{207}{96}nn\varepsilon
 -\tfrac94n^3\varepsilon\right)\sin a
 -\tfrac{27}{32}n\varepsilon\varepsilon\sin2a,
\]
\noindent und sonach wären noch folgende Theile hinzuzufügen:
\begin{center}
\makebox[.75\linewidth][s]{.\hfill.\hfill.\hfill.\hfill.\hfill.\hfill.\hfill.\hfill.\hfill.\hfill.\hfill.\hfill.\hfill.\hfill.\hfill.\hfill.\hfill.\hfill.\hfill.\hfill.\hfill.\hfill.\hfill.\hfill.\hfill.\hfill.\hfill.\hfill.\hfill.\hfill.\hfill.\hfill.\hfill.}\\[4pt]
\makebox[.75\linewidth][s]{.\hfill.\hfill.\hfill.\hfill.\hfill.\hfill.\hfill.\hfill.\hfill.\hfill.\hfill.\hfill.\hfill.\hfill.\hfill.\hfill.\hfill.\hfill.\hfill.\hfill.\hfill.\hfill.\hfill.\hfill.\hfill.\hfill.\hfill.\hfill.\hfill.\hfill.\hfill.\hfill.\hfill.}
\end{center}
\par\noindent\hspace*{\parindent}Auf diese Weise findet sich mit Weglassung aller Glieder der 5\textsuperscript{ten}
Ordnung, wo sie entweder allein stehen oder wo man ohnehin nicht auf sie
rechnen kann:
\[
\beta=g\times
\]
\begin{landscape}
\begingroup\small
\setlength{\tabcolsep}{3pt}
\renewcommand{\arraystretch}{1.15}
\begin{longtable}{r|p{7.8cm}r@{\quad}p{4.3cm}r}
\multicolumn{1}{c}{\(\sin\)}&\multicolumn{2}{c}{Nach unserer Form}&
\multicolumn{2}{c}{Nach \textsc{Mayers} Form}\\
\(\lambda\)&\(1-\tfrac38gg+\tfrac{15}{64}g^4-\tfrac{3}{32}nngg-\tfrac{33}{64}n^4\)&
&\(\text{Id.}-2nnee+\tfrac{25}{32}e^4\)&\\
\(3\lambda\)&\(-\tfrac{1}{24}gg\quad \mathord{.}\,\mathord{.}\,\mathord{.}\,\mathord{.}\,\mathord{.}\,\mathord{.}\,\mathord{.}\,\mathord{.}\,\mathord{.}\,\mathord{.}\,\mathord{.}\,\mathord{.}\,\mathord{.}\,\mathord{.}\,\mathord{.}\,\mathord{.}\,\mathord{.}\,\mathord{.}\)&\(-6\overset{\prime\prime}{,}28\)
&\(\text{Id.}\quad \mathord{.}\,\mathord{.}\,\mathord{.}\,\mathord{.}\,\mathord{.}\,\mathord{.}\,\mathord{.}\,\mathord{.}\,\mathord{.}\,\mathord{.}\)&\(-6\overset{\prime\prime}{,}28\)\\
\(2E-\lambda\)&\(\tfrac38n(1+\tfrac14n-\tfrac{3}{64}nn+\tfrac16\varepsilon\varepsilon-\tfrac12gg)\)&\(+528{,}67\)
&\(\text{Id.}\quad \mathord{.}\,\mathord{.}\,\mathord{.}\,\mathord{.}\,\mathord{.}\,\mathord{.}\,\mathord{.}\,\mathord{.}\,\mathord{.}\,\mathord{.}\)&\(+528{,}67\)\\
\(2E+\lambda\)&\(\tfrac{3}{32}ngg\quad \mathord{.}\,\mathord{.}\,\mathord{.}\,\mathord{.}\,\mathord{.}\,\mathord{.}\,\mathord{.}\,\mathord{.}\,\mathord{.}\,\mathord{.}\,\mathord{.}\,\mathord{.}\,\mathord{.}\,\mathord{.}\,\mathord{.}\,\mathord{.}\,\mathord{.}\,\mathord{.}\,\mathord{.}\)&\(+1{,}06\)
&\(\text{Id.}\quad \mathord{.}\,\mathord{.}\,\mathord{.}\,\mathord{.}\,\mathord{.}\,\mathord{.}\,\mathord{.}\,\mathord{.}\,\mathord{.}\,\mathord{.}\)&\(+1{,}06\)\\
\(2E-3\lambda\)&\(-\tfrac{3}{32}ngg\quad \mathord{.}\,\mathord{.}\,\mathord{.}\,\mathord{.}\,\mathord{.}\,\mathord{.}\,\mathord{.}\,\mathord{.}\,\mathord{.}\,\mathord{.}\,\mathord{.}\,\mathord{.}\,\mathord{.}\,\mathord{.}\,\mathord{.}\,\mathord{.}\,\mathord{.}\)&\(-1{,}06\)
&\(\text{Id.}\quad \mathord{.}\,\mathord{.}\,\mathord{.}\,\mathord{.}\,\mathord{.}\,\mathord{.}\,\mathord{.}\,\mathord{.}\,\mathord{.}\,\mathord{.}\)&\(-1{,}06\)\\
\(M+\lambda\)&\(\tfrac14nne-\tfrac{3}{32}n^3e\quad \mathord{.}\,\mathord{.}\,\mathord{.}\,\mathord{.}\,\mathord{.}\,\mathord{.}\,\mathord{.}\,\mathord{.}\,\mathord{.}\,\mathord{.}\,\mathord{.}\,\mathord{.}\,\mathord{.}\)&\(+1{,}39\)
&\(\text{Id.}\quad \mathord{.}\,\mathord{.}\,\mathord{.}\,\mathord{.}\,\mathord{.}\,\mathord{.}\,\mathord{.}\hspace{.65em}\mathord{.}\,\mathord{.}\)&\(+1{,}39\)\\
\(M-\lambda\)&\(\tfrac94nne+\tfrac{27}{8}n^3e\quad \mathord{.}\,\mathord{.}\,\mathord{.}\,\mathord{.}\,\mathord{.}\,\mathord{.}\,\mathord{.}\,\mathord{.}\,\mathord{.}\,\mathord{.}\,\mathord{.}\,\mathord{.}\,\mathord{.}\)&\(+14{,}32\)
&\(\text{Id.}+\tfrac54e^3\quad \mathord{.}\,\mathord{.}\,\mathord{.}\,\mathord{.}\,\mathord{.}\,\mathord{.}\)&\(+16{,}93\)
\end{longtable}
\SourceLeaf{639}
\begin{longtable}{r|p{7.8cm}r@{\quad}p{4.3cm}r}
\multicolumn{1}{c}{\(\sin\)}&&&&\\
\(2M-\lambda\)&\(\tfrac58ee\quad \mathord{.}\,\mathord{.}\,\mathord{.}\,\mathord{.}\,\mathord{.}\,\mathord{.}\,\mathord{.}\,\mathord{.}\,\mathord{.}\,\mathord{.}\,\mathord{.}\,\mathord{.}\)&\(+23\overset{\prime\prime}{,}68\)
&\(\text{Id.}-\tfrac94nnee\quad \mathord{.}\,\mathord{.}\,\mathord{.}\,\mathord{.}\,\mathord{.}\)&\(+22\overset{\prime\prime}{,}97\)\\
\(3M-\lambda\)&\(0\quad \mathord{.}\,\mathord{.}\,\mathord{.}\,\mathord{.}\,\mathord{.}\,\mathord{.}\,\mathord{.}\,\mathord{.}\,\mathord{.}\,\mathord{.}\,\mathord{.}\,\mathord{.}\,\mathord{.}\,\mathord{.}\,\mathord{.}\,\mathord{.}\,\mathord{.}\)&\(0\)
&\(-\tfrac54e^3\quad \mathord{.}\,\mathord{.}\,\mathord{.}\,\mathord{.}\,\mathord{.}\,\mathord{.}\,\mathord{.}\,\mathord{.}\,\mathord{.}\,\mathord{.}\)&\(-2{,}61\)\\
\(2E+M-\lambda\)&\(-\tfrac14nne-\tfrac{79}{96}n^3e\quad \mathord{.}\,\mathord{.}\)&\(-1{,}78\)
&\(\text{Id.}-\tfrac{75}{32}ne^3\quad \mathord{.}\,\mathord{.}\,\mathord{.}\,\mathord{.}\,\mathord{.}\)&\(-2{,}15\)\\
\(2E-M-\lambda\)&\(\tfrac94nne+\tfrac{45}{8}n^3e\quad \mathord{.}\,\mathord{.}\,\mathord{.}\,\mathord{.}\)&\(+15{,}28\)
&\(\text{Id.}\quad \mathord{.}\,\mathord{.}\,\mathord{.}\,\mathord{.}\,\mathord{.}\,\mathord{.}\,\mathord{.}\,\mathord{.}\,\mathord{.}\,\mathord{.}\,\mathord{.}\,\mathord{.}\)&\(+15{,}28\)\\
\(2E-2M-\lambda\)&\(-\tfrac{15}{16}nee-\tfrac{3}{64}nnee\)&\(-3{,}96\)
&\(\text{Id.}-\tfrac94nnee\quad \mathord{.}\,\mathord{.}\,\mathord{.}\,\mathord{.}\,\mathord{.}\)&\(-4{,}67\)\\
\(2E-2M+\lambda\)&\(\tfrac{15}{64}nee\quad \mathord{.}\,\mathord{.}\,\mathord{.}\,\mathord{.}\,\mathord{.}\,\mathord{.}\,\mathord{.}\,\mathord{.}\,\mathord{.}\,\mathord{.}\)&\(+1{,}05\)
&\(\text{Id.}+\tfrac{55}{64}nnee\quad \mathord{.}\,\mathord{.}\,\mathord{.}\,\mathord{.}\)&\(+1{,}28\)\\
\(\lambda+a\)&\(-\tfrac98n\varepsilon+\tfrac{69}{64}nn\varepsilon\quad \mathord{.}\,\mathord{.}\)&\(-24{,}37\)
&\(0\quad \mathord{.}\,\mathord{.}\,\mathord{.}\,\mathord{.}\,\mathord{.}\,\mathord{.}\,\mathord{.}\,\mathord{.}\,\mathord{.}\,\mathord{.}\,\mathord{.}\,\mathord{.}\,\mathord{.}\,\mathord{.}\,\mathord{.}\,\mathord{.}\,\mathord{.}\)&\(0\)\\
\(\lambda-a\)&\(\tfrac98n\varepsilon-\tfrac{69}{64}nn\varepsilon\quad \mathord{.}\,\mathord{.}\,\mathord{.}\,\mathord{.}\)&\(+26{,}01\)
&\(\tfrac{15}{16}nn\varepsilon\quad \mathord{.}\,\mathord{.}\,\mathord{.}\,\mathord{.}\,\mathord{.}\,\mathord{.}\,\mathord{.}\,\mathord{.}\,\mathord{.}\)&\(+1{,}64\)\\
\(\lambda+2a\)&\(\tfrac{27}{32}n\varepsilon\varepsilon\quad \mathord{.}\,\mathord{.}\,\mathord{.}\,\mathord{.}\,\mathord{.}\,\mathord{.}\,\mathord{.}\,\mathord{.}\,\mathord{.}\,\mathord{.}\)&\(+0{,}32\)
&\(0\quad \mathord{.}\,\mathord{.}\,\mathord{.}\,\mathord{.}\,\mathord{.}\,\mathord{.}\,\mathord{.}\,\mathord{.}\,\mathord{.}\,\mathord{.}\,\mathord{.}\,\mathord{.}\,\mathord{.}\,\mathord{.}\,\mathord{.}\,\mathord{.}\,\mathord{.}\)&\(0\)\\
\(\lambda-2a\)&\(-\tfrac{27}{32}n\varepsilon\varepsilon\quad \mathord{.}\,\mathord{.}\,\mathord{.}\,\mathord{.}\,\mathord{.}\,\mathord{.}\,\mathord{.}\,\mathord{.}\)&\(-0{,}32\)
&\(0\quad \mathord{.}\,\mathord{.}\,\mathord{.}\,\mathord{.}\,\mathord{.}\,\mathord{.}\,\mathord{.}\,\mathord{.}\,\mathord{.}\,\mathord{.}\,\mathord{.}\,\mathord{.}\,\mathord{.}\,\mathord{.}\,\mathord{.}\,\mathord{.}\,\mathord{.}\)&\(0\)\\
\(2E-\lambda+a\)&\(-\tfrac38n\varepsilon+\tfrac{45}{64}nn\varepsilon\quad \mathord{.}\,\mathord{.}\)&\(-7{,}53\)
&\(\text{Id.}-\tfrac{27}{64}nn\varepsilon\quad \mathord{.}\,\mathord{.}\,\mathord{.}\,\mathord{.}\,\mathord{.}\)&\(-8{,}26\)\\
\(2E-\lambda-a\)&\(-\tfrac18n\varepsilon-\tfrac{53}{64}nn\varepsilon\quad \mathord{.}\,\mathord{.}\)&\(-4{,}36\)
&\(\text{Id.}+\tfrac{27}{64}nn\varepsilon\quad \mathord{.}\,\mathord{.}\,\mathord{.}\,\mathord{.}\,\mathord{.}\)&\(-3{,}63\)\\
\(2E-\lambda+2a\)&\(\tfrac{15}{32}n\varepsilon\varepsilon\quad \mathord{.}\,\mathord{.}\,\mathord{.}\,\mathord{.}\,\mathord{.}\,\mathord{.}\,\mathord{.}\,\mathord{.}\,\mathord{.}\,\mathord{.}\)&\(+0{,}18\)
&\(\text{Id.}\quad \mathord{.}\,\mathord{.}\,\mathord{.}\,\mathord{.}\,\mathord{.}\,\mathord{.}\,\mathord{.}\,\mathord{.}\,\mathord{.}\,\mathord{.}\,\mathord{.}\,\mathord{.}\)&\(+0{,}18\)\\
\(2E-\lambda-2a\)&\(\tfrac18n\varepsilon\varepsilon\quad \mathord{.}\,\mathord{.}\,\mathord{.}\,\mathord{.}\,\mathord{.}\,\mathord{.}\,\mathord{.}\,\mathord{.}\,\mathord{.}\,\mathord{.}\)&\(+0{,}04\)
&\(\text{Id.}\quad \mathord{.}\,\mathord{.}\,\mathord{.}\,\mathord{.}\,\mathord{.}\,\mathord{.}\,\mathord{.}\,\mathord{.}\,\mathord{.}\,\mathord{.}\,\mathord{.}\,\mathord{.}\)&\(+0{,}04\)
\end{longtable}
\endgroup
\end{landscape}
\par\noindent\hspace*{\parindent}Diese Formel stimmt mit den Tafeln durchgängig so genau überein, als
man sich auf diese selbst verlassen kann.
\begin{center}
\rule{3cm}{.4pt}
\end{center}
\clearpage
\endgroup
